\section{Level 1.2: Initial Open Discovery}
\label{sec:level12}

\subsection{Experimental design}
\label{sec:level12-design}

Level 1.2's own Stage 0 preregistration manifest freezes the discovery
pipeline's commitments (three candidate clustering methods, $k=1$--$8$,
three missingness strategies), the twelve recurrence criteria as a
conjunction, and a mandatory mechanism-substream random-number-generator
mode for every causal run, all before any Stage 1--7 run executed
(\cref{tab:l12-design}).

\begin{table}[htbp]
\centering
\caption{Level 1.2 experimental design by stage.}
\label{tab:l12-design}
\adjustbox{max width=\textwidth}{\begin{tabular}{@{}lrrl@{}}
\toprule
\textbf{Stage} & \textbf{Seeds/draws} & \textbf{Total runs} & \textbf{Purpose} \\
\midrule
1: mechanism activation & 8 seeds $\times$ 7 conditions & 56 & Confirm every mechanism reachable \\
2: small open pilot      & 10 seeds, 1 condition        & 10 & First open-ended exploratory look \\
3: main open search       & 12 draws $\times$ 2 seeds     & 24 & 8-dim Latin Hypercube parameter sweep \\
4: blind discovery        & pools Stage 3               & (24 pooled) & Clustering + 12 recurrence criteria \\
5: prospective confirmation & 12 draws (fresh seeds)     & 12 & Independent replication of Stage 4 \\
6: scaling robustness      & 4 conditions $\times$ 2 seeds & 8 & Population/patch-count scale grid \\
7: perturbation/transition & 2 types $\times$ 3 seeds      & 6 pairs & Shock recovery/transition test \\
13: causal ablations       & 2 factors $\times$ 8 matched pairs & 32 & Matched-seed causal test \\
\bottomrule
\end{tabular}
}
\end{table}

Seven stages followed: mechanism-activation
calibration (Stage 1); a small, single-condition open pilot (Stage 2, $n=10$
seeds); a main open evolutionary search over an eight-dimensional Latin
Hypercube parameter space (Stage 3, 12 parameter draws $\times$ 2 seeds =
24 runs); a blind discovery run over Stage 3's pooled output (Stage 4); an
independent, fresh-seed prospective confirmation using the same 12
parameter draws (Stage 5, $n=12$); a scaling-robustness grid (Stage 6,
population $\in\{40,80\}\times$ patches $\in\{3,6\}$, 8 runs); and a
perturbation/transition test (Stage 7, 6 shocked/unshocked pairs). Section
13 added two matched-seed causal ablations on top of Stages 2--7's purely
correlational findings.

\subsection{Collapse-dominated search}
\label{sec:level12-collapse}

The dominant finding across every stage was near-universal ecological
collapse. Stage 2's ten-seed pilot returned 7 collapse, 3 fragmented, 0
viable. Stage 3's 24-run main search returned 22 collapse, 2 fragmented, 0
viable --- collapse dominated across the swept parameter range, not one
isolated setting (\cref{fig:l12-stage3-outcomes}). Stage 5's fresh-seed
confirmation, reusing Stage 3's own 12 parameter draws, returned 11
collapse and, for the first time in the milestone, 1 viable outcome (draw
9, notably high patch regeneration rate). Stage 6's scaling-robustness
grid returned 8 collapse out of 8 runs across population 40--80 and
3--6 patches (\cref{fig:l12-scaling}) --- ruling out a small-population
artifact as the explanation. Stage 7's perturbation/transition test was
inconclusive by design: both shocked and unshocked baseline conditions had
already collapsed, leaving no non-collapse state for a shock to transition
away from, a disclosed floor effect rather than evidence that
perturbations are inert.

Section 13's two matched-seed causal ablations (\texttt{rng\_mode=%
"mechanism\_substreams"}, 8 seed pairs each) tested the two candidate
explanations Stages 2--7 had only observed correlationally. Toggling
migration on versus off at matched seeds found 5 of 8 pairs favored
migration off on ecological resource fraction (mean paired difference
$-0.077$), with 4 of 8 pairs showing an outcome-label transition (3 of
those 4 moving from \texttt{fragmented} to \texttt{collapse} when migration
was turned on); a Wilcoxon signed-rank test on the paired differences gave
$p=0.25$ \parencite{wilcoxon1945individual}, reported as supplementary
only given the small sample. Toggling the patch regeneration rate between
Stage 3's own sampled bounds (0.2 versus 1.0) found a more consistent
signal: 7 of 8 pairs favored the higher regeneration rate (mean paired
difference $+0.033$), 2 of 8 outcome-label transitions (both
\texttt{collapse}$\to$\texttt{fragmented}), Wilcoxon $p=0.039$. Neither
ablation overturned the headline collapse finding: 13 of 16
migration-ablation runs and 14 of 16 regeneration-ablation runs still
landed in \texttt{collapse}. Both factors are real, causally implicated
contributors; neither is sufficient alone to make \texttt{viable} the
typical outcome at the tested population and duration scale.

\subsection{Null recurrence result}
\label{sec:level12-null}

Stage 4's blind discovery run pooled Stage 3's 24 seed-run phenotype
vectors and evaluated 76 candidate clusters (two clustering methods across
8 values of $k$, plus \texttt{HDBSCAN}'s own clusters) against all twelve
recurrence criteria. Zero of 76 met all twelve
(\cref{fig:l12-recurrence-criteria}); the thirteen best candidates, each
passing eleven of twelve, were uniformly \texttt{hierarchical} $k=6$
partitions dominated by collapse-labeled members, and uniformly failed
\texttt{cross\_algorithm\_support} --- no other clustering method
corroborated that specific partition. Stage 5's independent replication,
on a genuinely fresh, disjoint 12-run pool using the identical, unmodified
discovery pipeline, found the same result: zero of 72 candidates met all
twelve criteria (\cref{fig:l12-discovery-replication}). Held-out
log-likelihood, computed by fitting on Stage 3's own training data and
scoring Stage 5's genuinely held-out fresh runs, degraded gradually from
$k=1$ to $k=3$ and then collapsed catastrophically at $k\geq4$ --- the
textbook signature of small-sample Gaussian-mixture overfitting, and
independent evidence against trusting any higher-$k$ partition from a
sample this size, consistent with the recurrence-criteria result.

Per this milestone's own closing instruction, this is reported exactly as
"the recurrent-phenotype hypothesis was not supported at this sample
size" --- not as evidence against any particular count, not zero, not
one, not three, and not as grounds to alter the model or widen the search
now that the outcome was known.

\subsection{Why ecological confounding required Level 1.2b}
\label{sec:level12-why-1.2b}

Level 1.2's own main pool put 22 to 24 of every 24 runs in
\texttt{collapse}. This licenses a narrower conclusion than "no recurrent
institutional configuration exists in this model": it licenses "within the
sampled region, collapse was the dominant recurrent \textit{system}
outcome, and no recurrent \textit{institutional} configuration was
independently supported" --- because a sample this collapse-saturated
gives a discovery pipeline very little non-collapse institutional
variation to find structure in even if such structure exists. Level
1.2's own completion report states this limitation explicitly: the
eligible, non-collapse subset available to any future discovery attempt
was 2 of 24 runs in Stage 3 and 1 of 12 in Stage 5. Level 1.2b
(\cref{sec:level12b}) exists to close exactly this inferential gap ---
not by re-running Level 1.2's search hoping for a different answer, but
by building an ex-ante ecological feasibility measure, stratifying the
design by it before any institutional variation is sampled, and sampling
at roughly 24 times the scale.
